This appendix explains how to run the submitted Iteration 3 Python/Jupyter
workflow. All commands below are run from the \path{Iteration3_OpenStack}
directory unless stated otherwise. Keep the supplied directory structure intact.

\subsection{Environment and Required Files}

Install Python and create a dedicated virtual environment. The submitted
notebooks record Python 3.13.5; package requirements are minimum versions rather
than an exact environment lock. On macOS or Linux, run:

\begin{verbatim}
cd INFOSYS722/Iteration3_OpenStack
python3 -m venv .venv
source .venv/bin/activate
python -m pip install -r requirements.txt
python -m ipykernel install --sys-prefix \
  --name gallstone-iteration3 \
  --display-name "Gallstone Iteration 3"
\end{verbatim}

On Windows PowerShell, create the environment with
\texttt{python -m venv .venv} and activate it using
\path{.venv\Scripts\Activate.ps1}. Run the kernel-install command on one line,
omitting the shell continuation backslashes.

The active modelling workflow requires \path{data/raw/dataset_gallstone},
which is a CSV file with no extension. Preserve its filename and row order.
The accompanying \path{data/raw/dataset_gallstone.xlsx} is also supplied, but
does not replace the extensionless CSV used by the modelling functions.
Keep \path{utils.py}, \path{modeling.py}, \path{requirements.txt},
\path{run_workflow.py}, and the \path{notebooks/}, \path{scripts/},
\path{tests/} and \path{Report/} directories together with the raw data.

\subsection{Run the Complete Workflow}

The runner executes the stages sequentially and stops if a stage fails:

\begin{verbatim}
python run_workflow.py --list
python run_workflow.py
\end{verbatim}

The first command only lists the stages; the second runs the complete analysis,
regenerates report figures and tables, and runs the workflow checks. Execution
overwrites notebook outputs, generated data, result files and saved models.
Use a separate copy of the submission to preserve the supplied results.
Each execution records Python and installed package versions in
\path{results/runs/}. Leave-one-out evaluation and stacking involve many model
fits and can take substantially longer than the other stages.

The execution order and main outputs are:

\begin{enumerate}[leftmargin=*]
  \item \textbf{Exploration (\texttt{eda}):}
  \path{notebooks/gallstone_eda.ipynb} explores the original data.
  \item \textbf{Preparation (\texttt{prepare}):}
  \path{notebooks/transformation_feature_engineering.ipynb} exports
  training-fitted transformations and partitions to \path{data/processed/}.
  \item \textbf{Holdout evaluation (\texttt{holdout}):}
  \path{notebooks/holdout_method.ipynb} writes metrics and predictions to
  \path{results/holdout/} and pipelines to \path{models/holdout/}.
  \item \textbf{Leave-one-out evaluation (\texttt{lopo}):}
  \path{notebooks/leave_one_out_method.ipynb} evaluates 319 held-out rows
  per model family and writes \path{results/lopo/} and \path{models/lopo/}.
  \item \textbf{Parameter tuning (\texttt{tuning}):}
  \path{notebooks/tuning_model.ipynb} writes search and evaluation results
  to \path{results/tuning/} and selected models to \path{models/tuned/}.
  \item \textbf{Error analysis (\texttt{errors}):}
  \path{notebooks/error_analysis.ipynb} reads saved predictions and writes
  \path{results/error_analysis/}.
  \item \textbf{Iteration comparison (\texttt{patterns}):}
  \path{scripts/iteration_pattern_analysis.py} compares the 20- and
  10-predictor references and writes \path{results/iteration_patterns/}.
  \item \textbf{Report outputs (\texttt{figures}):}
  \path{Report/build_modeling_figures.py} regenerates modelling figures
  and evidence tables in \path{Report/figures/modeling/} and
  \path{Report/generated/} from saved results.
  \item \textbf{Verification (\texttt{check}):} unit checks cover preparation
  isolation, saved-pipeline prediction consistency and disease-positive
  label/false-negative-rate coding.
\end{enumerate}

For interactive execution, run \texttt{python -m jupyter lab}, select the
\textbf{Gallstone Iteration 3} kernel and use \textbf{Restart Kernel and Run
All Cells} for each of the six notebooks in the order above. Complete each
notebook before opening the next stage, then run:

\begin{verbatim}
python scripts/iteration_pattern_analysis.py
python Report/build_modeling_figures.py
python -m unittest discover -s tests -v
\end{verbatim}

\subsection{Resume Execution and Check Results}

To restart at a failed stage, use \texttt{--from-stage}. To run just one stage,
also specify \texttt{--to-stage}, for example:

\begin{verbatim}
python run_workflow.py --from-stage tuning
python run_workflow.py --from-stage holdout --to-stage holdout
python run_workflow.py --from-stage check
\end{verbatim}

Stage selection does not rebuild prerequisites: tuning requires prepared data,
and complete error analysis requires the preceding evaluation results. After
changing data, seeds, preparation or model settings, rerun from
\texttt{prepare} through the end; rerun exploration when raw data change.
An interrupted stage is not a completed result.

Check \path{data/processed/manifest.json} and split-membership files for seeds
42, 7, 21, 84 and 123, with 223 training, 48 validation and 48 test rows per
seed. A complete fixed holdout run produces 120 metric rows and 5,760 prediction
rows; LOPO produces six metric rows and 1,914 predictions; tuning produces
25 metric rows and 1,200 predictions. These appear in the respective
\path{results/holdout/}, \path{results/lopo/} and \path{results/tuning/}
directories as \path{metrics.csv} and \path{predictions.csv}.
The checks verify selected workflow invariants rather than every reported score.

Fixed seeds support repeatability within an environment, but package and
platform differences can change exact results. Supplied pickle files require
compatible library versions; if incompatible, regenerate preparation and
models from \texttt{prepare} onward in the new environment. Evaluation must
use the training-fitted partitions, not a new random split of the full-data
\path{transformed_dataset.csv}.

\subsection{Rebuild the Report PDF}

With saved results present, regenerate the modelling figures and tables without
retraining using the first command below. Install a LaTeX distribution with
\texttt{latexmk}, \texttt{pdflatex}, Biber and the packages declared in
\path{Report/main.tex}, then compile:

\begin{verbatim}
python Report/build_modeling_figures.py
latexmk -pdf -interaction=nonstopmode -halt-on-error \
  -cd Report/main.tex
\end{verbatim}

On PowerShell, enter the \texttt{latexmk} command on one line without the
continuation backslash. The compiled report is \path{Report/main.pdf}.
Retain all report source files, \path{references.bib}, \path{media/}, figures
and generated tables. Figure generation does not update earlier exploratory
images or manually written tables and discussion; reconcile those passages
with any changed results before compiling the final report.
